
Transformer attention scales quadratically with sequence length, motivating research into efficient variants. Linear attention approximates softmax with kernel methods. Sparse attention restricts the attention pattern. Iterative refinement approaches propose that attention weights evolve across internal steps to converge on task-relevant patterns. These methods report compelling efficiency numbers, but the connection between proposed mechanism and observed efficiency is typically assumed rather than tested.

End-to-end benchmarks conflate mechanism with outcome. When temporal dynamic attention achieves FLOP reduction, does the reduction stem from attention patterns converging toward stable distributions, or from simpler architectural factors such as reduced iteration count? Without decomposing efficiency claims into testable sub-hypotheses, these possibilities cannot be distinguished.

This gap matters because mechanistic understanding determines generalizability. If efficiency derives from convergence, scaling temporal steps should amplify the benefit. If efficiency derives from iteration count, the benefit is bounded regardless of sequence length. These predictions diverge sharply, yet typical evaluation cannot distinguish them.

This work decomposes temporal dynamic attention efficiency claims into four sub-hypotheses with explicit gates and falsification criteria:

\begin{enumerate}
\item \textbf{H-E1 (Existence):} The mechanism can be implemented and trains stably.
\item \textbf{H-M1 (Efficiency):} Reducing temporal steps decreases FLOPs while maintaining perplexity.
\item \textbf{H-M2 (Convergence):} Attention patterns converge (entropy decreases) across temporal steps.
\item \textbf{H-C1 (Scaling):} Efficiency gains scale with sequence length.
\end{enumerate}

Results show that H-E1 and H-M1 pass while H-M2 and H-C1 fail. The efficiency gain derives from reduced iteration count, not attention convergence.

